\documentclass[10pt,a4paper]{amsart}
\usepackage[margin=19mm]{geometry}
\usepackage{amsmath,amssymb,mathtools}
\usepackage[scr=rsfs,cal=euler]{mathalfa}
\usepackage{xfrac}
\usepackage[unicode,hidelinks,bookmarks=false]{hyperref}
\newcommand{\ie}{i.e.\ }
\newcommand{\cf}{cf.\ }
\newcommand{\resp}{respectively\ }
\newcommand{\Subsection}{\subsection}
\providecommand{\nobreakdash}{\nobreak\hbox{-}}
\setlength{\emergencystretch}{2em}
\multlinegap0pt
\usepackage{bm}
\usepackage{bbm}
\usepackage{dsfont}
\usepackage{xspace}
\usepackage{cancel}
\usepackage{mathtools}
\usepackage{amsthm}
\makeatletter
\@ifundefined{thm}{\newtheorem{thm}{Theorem}}{}
\@ifundefined{proposition}{\newtheorem{proposition}[thm]{Proposition}}{}
\makeatother
\title[The largest fragment in self-similar fragmentation processes]{The largest fragment in self-similar fragmentation processes of positive index}
\author{Piotr Dyszewski, Samuel G. G. Johnston, Sandra Palau, Joscha Prochno}
\date{Source: arXiv:2409.11795v3 (CC BY 4.0)}
\begin{document}
\maketitle

\noindent\emph{Source and licence.} The largest fragment in self-similar fragmentation processes of positive index. Version arXiv:2409.11795v3 (CC BY 4.0), distributed under the Creative Commons Attribution 4.0 International licence, as recorded in the arXiv licence metadata for this exact version and shown on the abstract page https://arxiv.org/abs/2409.11795v3. Full rights notice: licenses/arxiv-2409.11795-CC-BY-4.0.txt.\par\smallskip

\noindent\textit{Editorial scope.} Counted unit: the Thm whose source label is \texttt{thm:antichains} in the article (source0.tex lines 1228--1237, source label \texttt{thm:antichains}), delivered together with the article's own complete original proof of it (source0.tex lines 1239--1308, one \texttt{proof} environment of 70 lines). No mathematics is written, repaired or re-derived: statement and proof are the article's own text line for line. Required context retained uncounted: prop:heightlower (lines 1173--1178). Every definition, display and label of those lines is retained unchanged. Omitted from this package and not redistributed: 1--1172, 1179--1227, 1238--1238, 1309--3177. Nothing inside a retained excerpt is omitted or altered: every retained body is delivered exactly as the article's lines are written, so no omission receipt of any kind accompanies this package. Citations: the retained excerpts carry no citation command at all, so no bibliography entry of the article is transcribed and no reference is redistributed. Reference adaptation: none was needed and none was made; every reference inside the delivered unit resolves inside it, so no unresolved reference or citation is produced. Printed slips kept literal: no printed word, symbol, hypothesis or number of the counted statement or of its proof was altered. Layout and class: the article's own class is \texttt{amsart} with packages including inputenc, amsfonts, amssymb, amsmath, comment, xcolor, hyperref, graphicx, sidecap, accents, tikz, geometry; the wrapper is the lane's pinned \texttt{amsart} editorial wrapper at 10pt on A4, which restates the theorem-like environments the retained text needs and adds the article's own macro definitions that the retained text needs (none), with extra wrapper packages (bm, bbm, dsfont, xspace, cancel, mathtools). The delivered numbering is the unit's own. Assets: no retained span contains a figure environment, a graphics-inclusion command, a PDF-page inclusion, a table or a TikZ picture; no image file is copied into this package, the package declares no asset dependency and no image file is redistributed with it. Licence: version arXiv:2409.11795v3 of this article is distributed under the Creative Commons Attribution 4.0 International licence, as recorded in the arXiv licence metadata for this exact version and shown on the abstract page https://arxiv.org/abs/2409.11795v3. A full rights notice is delivered at licenses/arxiv-2409.11795-CC-BY-4.0.txt. Characters: every retained excerpt is pure ASCII, so the delivered engine, XeLaTeX with its default Latin Modern text font, reports no missing character.
\begin{proposition} \label{prop:heightlower}
Let $a > 0$ and $\varepsilon > 0$. Then, there exists a random number $h(a,\varepsilon)$ such that for all $h \geq h(a,\varepsilon)$ we have
\begin{align*}
D_h^a \geq e^{ (1 - \varepsilon) h}.
\end{align*}
\end{proposition}

	\begin{thm} \label{thm:antichains}
		Let $\delta > 0 $ and  $ \log(2) > a > 0$. 
		Then there exists a random number $h_0(a,\delta)\in[0,\infty)$ such that for all $h \geq h_0(a,\delta)$, 
		there exists a subset $\mathbb{S}$ of $\mathbb{T}$ so that the following hold:
		\begin{enumerate}
			\item $\mathbb{S}$ is an antichain,
			\item $\mathcal{H}(w) \in (h-a,h]$ for all $w \in \mathbb{S}$,
			\item The cardinality of $\mathbb{S}$ satisfies $\# \mathbb{S} \geq e^{ (1 - \delta)h}$.
		\end{enumerate}
	\end{thm}

	\begin{proof}
		Define
		\begin{align*}
			\mathcal{D}(h-a,h] := \{ w \in \mathbb{T} : \mathcal{H}(w) \in (h-a,h] \}.
		\end{align*}
        
		By Proposition \ref{prop:heightlower} there exists a random variable 
		$h_0(a,\delta/2)$ such that for all $h \geq h_0(a,\delta/2)$ we have
		\begin{align} \label{eq:bc1}
			D_h^a = \# \mathcal{D}(h-a,h] \geq e^{ (1 - \delta/2) h}. 
		\end{align}
		We now show that with high probability, $\mathcal{D}(h-a,h]$ contains 
		an antichain of size at least $e^{(1-\delta)h}$.
		To this end, let $\tilde{\mathcal{D}}(h-a,h]$ denote the set of first ancestors in $(h-a,h]$. 
		This is the subset of $\mathcal{D}(h-a,h]$ consisting of those elements with no 
	ancestor also in $\mathcal{D}(h-a,h]$. In other words,
		\begin{align*}
			\tilde{\mathcal{D}}(h-a,h] := 
			\{w \in \mathcal{D}(h-a,h] : \mathcal{H}(v) \leq h-a ~\text{for every strict ancestor $v$ of $w$} \}.
		\end{align*}
		Plainly, $\tilde{\mathcal{D}}(h-a,h]$ is an antichain. 
		We now seek to bound below the cardinality $\tilde{D}_h^a$ of $\tilde{\mathcal{D}}(h-a,h]$. 
		In this direction, let $w \in \tilde{\mathcal{D}}(h-a,h]$ and define the random variable
		\begin{align*}
			N_w := \# \{ v \in \mathcal{D}(h-a,h] : v \text{ descendant of $w$} \}.
		\end{align*}
		We control the expectation of $N_w$. Recall that from the conservativeness of the underlying fragmentation process for every $w \in \mathbb{T}$ we have
		\begin{align*} 
			1 = \sum_{i \geq 1}e^{ - ( \mathcal{H}(wi) - \mathcal{H}(w) ) }.
		\end{align*}
		Note that since $a < \log(2)$, the previous equation tells us that each $w$ in $\mathbb{T}$ 
		may have at most one child $wi$ with $\mathcal{H}(wi) - \mathcal{H}(w) < a$. 
		Let $p(a) \in (0,1)$ denote the probability of this occurring, i.e.,
		\begin{align*}
			p(a) := \mathbb{P}\left[ \exists i : \mathcal{H}(i)  < a \right].
		\end{align*}
		It follows that for each $w$ in $\tilde{\mathcal{D}}(h-a,h]$, $N_w$ is stochastically dominated by a 
		geometric random variable with success probability $p(a)$. In particular, there is some constant $c(a) < \infty$ such that 
		\begin{align*}
			\mathbb{E}\big[N_w+1 | w \in \mathcal{D}(h-a,h]\big] \leq c(a).
		\end{align*}
		In particular, 
		\begin{align*}
			\mathbb{E}\left[ D_h^a \left| \tilde{D}_h^a\right. \right] \leq c(a) \tilde{D}_h^a.
		\end{align*}
		By a conditional Markov's inequality, it follows that
		\begin{align} \label{eq:bc2}
			\mathbb{P}\left[ \tilde{D}_h^a \leq e^{ -\delta h/2} D_h^a \right] \leq c(a) e^{ - \delta h/2}.
		\end{align}
		Combining \eqref{eq:bc1} with \eqref{eq:bc2} and using the Borel--Cantelli lemma, 
		it follows that there exists a random $h_1(a,\delta)$ such that 
		\begin{align*}
			\# \tilde{\mathcal{D}}(h-a,h] \geq e^{(1-\delta)h} \qquad 
			\text{for all $h \geq h_1(a,\delta)$ such that $ h  = na $ for some $n \in \mathbb{N}$}.  
		\end{align*} 
		To transfer this result from $h$ sufficiently large taking values in the lattice 
		$\{ n a : n \in \mathbb{N} \}$ to all sufficiently large $h$, 
		note that any interval of the form $(h-a,h]$ necessarily contains a subinterval of the form $((n-1)a/2,na/2]$. In particular,
		\begin{align*}
			\tilde{\mathcal{D}}(h-a,h] \supseteq \tilde{\mathcal{D}}\left(n\frac{a}{2} - \frac{a}{2}, n\frac{a}{2} \right] 
			\qquad \text{for some $n \in \mathbb{N}$}.
		\end{align*}
		It thus follows that for $h \geq h_2(a,\delta) := h_1(a/2,\delta)$ we have
		\begin{align*}
			\# \tilde{\mathcal{D}}(h-a,h] \geq e^{(1-\delta)h}.
		\end{align*}
		Thus, for all real $h$ greater than some random height $h_2(a,\delta)$, 
		the antichain $\mathbb{S} := \tilde{\mathcal{D}}(h-a,h]$ of first ancestors in $(h-a,h]$ 
		has cardinality at least $e^{(1-\delta)h}$, completing the proof.
\end{proof}

\end{document}
